% !PRINTED 304
\weberpage{304}

%<<CONT
\gesperrt{annehmen, und $\varrho$ durchläuft, je nachdem $\alpha^{-1}\beta$ quadratischer Rest oder Nichtrest ist, in einem System die Reihe der Reste oder der Nichtreste.} Die Änderung des Vorzeichens von $\xi$ gibt keine neue Substitution $W$. Hiernach können wir $W$ so darstellen:

\weq{21}{W=\begin{pmatrix}\xi, & \xi\varrho\\ \xi^{-1}\varrho^{-1}(m-1), & \xi^{-1}m\end{pmatrix}.}

Da nun das Quadrat von $W$

\[ W^2=\begin{pmatrix}\xi^2+m-1, & (\xi^2+m)\varrho\\ \xi^{-2}\varrho^{-1}(m-1)(\xi^2+m), & [(m-1)\xi^2+m^2]\xi^{-2}\end{pmatrix} \]

zu $\mathfrak{K}$ gehören muß, so ist es unter einem der Systeme (13), (14), (16), (17) enthalten, woraus sich folgende vier Möglichkeiten ergeben:

\weq{22}{\begin{aligned}
1.&\quad \xi^2+m\equiv 0,\\
2.&\quad \xi^2+m-1\equiv 0,\quad (m-1)\xi^2+m^2\equiv 0,\\
3.&\quad (\xi^2+m-1)[(m-1)\xi^2+m^2]\equiv m\xi^2,\\
4.&\quad (\xi^2+m-1)[(m-1)\xi^2+m^2]\equiv m'\xi^2,
\end{aligned}}

und jeder der $p-1$ Werte (20), für $\xi$ gesetzt, muß einem dieser vier Fälle genügen. Wenn eine der Kongruenzen (22), 2. erfüllt ist, so muß die andere daraus folgen, was nur möglich ist, wenn

\weq{23}{2m\equiv 1,\qquad \xi^2\equiv m\quad(\mathrm{mod}\ p).}

Nun hat eine Kongruenz in bezug auf einen Primzahlmodul \gesperrt{höchstens} so viele inkongruente Wurzeln, als der Grad der Kongruenz beträgt; 1. und 2. können also höchstens für je zwei, 3. und 4. höchstens für je vier Werte von $\xi$ befriedigt sein; also gibt es im ganzen höchstens zwölf Werte von $\xi$, die \gesperrt{einer} der vier Kongruenzen (22) genügen.

Daraus folgt:

\[ p-1\leqq 12,\qquad p\leqq 13. \]

Ist aber $p=13$, so muß jede der Kongruenzen (22) die Maximalzahl von Wurzeln haben; es muß also auch 2. für zwei Werte von $\xi$ befriedigt sein; dann müßte nach (23) $m\equiv 7$, $\xi^2\equiv 7\ (\mathrm{mod}\ 13)$ sein, was unmöglich ist, da $7$ quadratischer Nichtrest von $13$ ist.

\gesperrt{Ein Teiler von $\mathfrak{L}_0$ vom Index $p$ existiert also nicht,} wenn $p>11$ ist.
